Wie werden die Abschlusskonten/Eröffnungskonten erstellt und fließen in die Bilanz ein

Gemäß dem Prinzip der doppelten Buchführen sollten auf jedem Konto sich die Aktiva/Passiva ausgleichen 
um die Information des Endbestandes auf einem Konto zu vermerken  und das Gleichgewicht aufrecht zu erhalten 
wird dieser nach der Saldierung (Endbestandes = Zugang - Abgang) auf der Haben/Soll Seite bei Aktiv/Passivkonten
jeweils eingetragen.

Als Beispiel hier das Aktivkonto Kasse

IMAGES FROM PAGE 27

Für die Bilanz wird dann nochmal ein separates Schlussbilanzkonto geführt um das korrespondierende 
Gegenkonto (Kontenruf) zum Endbestand Eintrag zu haben (jeder Buchungssatz muss einmal auf der Haben und einmal auf der Soll
erscheinen).

Es ergibt sich folgender Buchungssatz und Konten


IMAGES FROM PAGE 28

Für Passivkonten gilt dieselbe Methodik aber umgekehrt, hier ein verkürztes Beispiel:

IMAGES FROM PAGE 29

Das Eröffnen der Konten zu Anfang eines neuen Geschäftsjahres geschieht durch eine Buchung
in den Soll oben links und in den Haben Rechts bei Aktiv/Passiv Konten. Was sollte aber der
entsprechende Kontenruf sein? Die Lösung ist das Eröffnungsbilanzkonto:


wodurch zusammengesetzte Buchungssätze entstehen schließlich stimmt die Anzahl an Posten nicht mit den Einträgen
überein (jeduch mit den Geldbeträgen worauf es eigentlich ankommt).


Erklären sie kurz das Eigenkapitalkonto 

Das Eigenkapitalkonto ist notwendig zur Führung der GuV, da sich
die Höhe des Eigenkapitals aus dem Saldo von Vermögen und Fremdkapital ergibt.

Grundsätzlich könnte man also den Gewinn/Verlust (einer Periode) einfach nur erhalten
durch vergleich des Eigenkapitals zur Anfangs und Schlussbilanz was auch Distanzrechnung genannt wird.

INSERT IMAGE PAGE 34


Wie wird die Gewinn- und Verlustrechnung (GuV) nach § 242 II und IV HGB erstellt?

Aufwendungen sind definiert als periodisierte Ausgaben, d.h. als Ausgaben, die
der Periode zugerechnet werden, und Erträge analog als periodisierte Einnahmen.

Diese werden in Aufwands- und Ertragskonten geführt dessen Gegenkonto ist das 
Gewinn- und Verlustkonto dessen Saldo nach direkt bevor Abschluss in das Eigenkapitalkonto 
eingetragen wird. Die GuV wird repräsentiert durch die Gewinn- und Verlustkonten.


Beispiele für Aufwandsarten, für die eigene Konten eingerichtet werden, sind Lohn- und Gehalts-
aufwendungen, Materialverbräuche, Abschreibungen, Zinsaufwendungen, Auf-
wendungen für Fremdleistungen und Steueraufwendungen.

Typische Ertragsarten sind Umsatzerlöse, Zinserträge, Mieterträge, Beteiligungserträge 
aber auch beispielsweise staatliche Subventionen.

Ein Beispiel für die Übertragung von Ertragskonten bis zum Eigenkapitalkonto

INSERT IMAGE FROM PAGE 41

Für Aufwandskonten geschieht das analog wobei der Saldo im Haben startet.
WICHTIG ist das dass GUV Konto nicht notwendigerweise sein Saldo im Haben/Soll (wie im Bild) stehen 
hat, dies ist davon abhängig ob jeweils ein Jahresverlust oder Gewinn vorliegt.

Zu beachten sind das periodenfremden Aufwendungen und Erträge nicht miteinflüssen dürfen in die GuV
ebenso wie Außergewöhnliche Aufwendungen und Erträge (e.g. Brandschaden), beide sind stattdessen im Anhang.


Eine Besonderheit der Erfolgskonten gegenüber den Bestandskonten ist, dass die
Erfolgskonten keine Anfangsbestände haben. Erträge und Aufwendungen sind
Stromgrößen.
